\documentclass[12pt,a4paper]{article}
\usepackage[utf8]{inputenc}
\usepackage[spanish]{babel}
\usepackage{amsmath}
\usepackage{amssymb}
\usepackage{geometry}
\usepackage{array}
\usepackage{booktabs}
\usepackage{graphicx}
\usepackage{float}
\usepackage{multicol}
\usepackage{fancyhdr}
\usepackage{enumitem}
\usepackage{titlesec}
\usepackage{xcolor}
\usepackage{mdframed}
\usepackage{tcolorbox}
\usepackage{colortbl}

\usepackage{microtype}
\usepackage{hyperref}
\usepackage{pgfplots}
\pgfplotsset{compat=1.18}
\usepackage{tikz}

\tcbuselibrary{skins, breakable, theorems}

\geometry{margin=2.2cm, top=2.5cm, bottom=2.5cm}
\setlength{\headheight}{24.01996pt}
\addtolength{\topmargin}{-12pt}

% Paleta de colores
\definecolor{azulOscuro}{RGB}{25, 60, 110}
\definecolor{azulMedio}{RGB}{52, 101, 164}
\definecolor{azulClaro}{RGB}{214, 229, 248}
\definecolor{dorado}{RGB}{180, 140, 40}
\definecolor{grisClaro}{RGB}{245, 246, 250}
\definecolor{grisBorde}{RGB}{200, 210, 225}
\definecolor{rojoGrafico}{RGB}{180, 30, 30}
\definecolor{verdeGrafico}{RGB}{30, 130, 60}

% Encabezado
\pagestyle{fancy}
\fancyhf{}
\renewcommand{\headrulewidth}{0.6pt}
\renewcommand{\headrule}{\hbox to\headwidth{\color{azulMedio}\leaders\hrule height \headrulewidth\hfill}}
\rhead{\color{azulOscuro}\small\thepage}
\lhead{\color{azulOscuro}\small\textit{Formulario Unidad II --- Trigonometría}}

% Títulos
\titleformat{\section}
  {\color{azulOscuro}\large\bfseries}
  {\color{dorado}\thesection.}{0.5em}{}
  [\vspace{-4pt}\textcolor{azulMedio}{\rule{\linewidth}{0.5pt}}\vspace{2pt}]

\titleformat{\subsection}
  {\color{azulMedio}\normalsize\bfseries}
  {\thesubsection.}{0.5em}{}

% Cajas
\tcbset{
  formulabox/.style={
    enhanced, breakable,
    colback=grisClaro, colframe=azulMedio,
    fonttitle=\bfseries\color{white},
    coltitle=white, colbacktitle=azulOscuro,
    boxrule=0.6pt, arc=4pt,
    left=6pt, right=6pt, top=4pt, bottom=4pt,
  },
  notabox/.style={
    enhanced,
    colback=azulClaro!40, colframe=azulMedio,
    boxrule=0.5pt, arc=3pt,
    left=6pt, right=6pt, top=3pt, bottom=3pt,
    fontupper=\small,
  },
  ejercbox/.style={
    enhanced, breakable,
    colback=white, colframe=dorado,
    boxrule=0.7pt, arc=4pt,
    left=6pt, right=6pt, top=4pt, bottom=4pt,
  }
}

\begin{document}

% ── PORTADA ──────────────────────────────────────────────────────
\begin{titlepage}
  \pagecolor{azulOscuro}
  \centering
  \vspace*{3cm}
  {\color{white}\fontsize{13}{16}\selectfont\scshape
    Cálculo I --- Ingeniería Civil e Informática\par}
  \vspace{0.4cm}
  {\color{grisBorde}\rule{0.6\linewidth}{0.4pt}\par}
  \vspace{0.8cm}
  {\color{white}\fontsize{32}{38}\selectfont\bfseries Formulario\par}
  {\color{dorado}\fontsize{28}{34}\selectfont\bfseries Unidad II\par}
  \vspace{0.5cm}
  {\color{azulClaro}\fontsize{16}{20}\selectfont Trigonometría\par}
  \vspace{0.3cm}
  {\color{grisBorde}\small
    Identidades Trigonométricas $\cdot$
    Teorema del Seno y Coseno $\cdot$ Aplicaciones\par}
  \vspace{1.2cm}
  {\color{grisBorde}\small
    Preparado por : Diego Solís S1 $\cdot$
    Joaquin Vasquez S2 $\cdot$ Benjamin Romo S3 \par}
  \vspace{1.2cm}
  {\color{grisBorde}\rule{0.4\linewidth}{0.3pt}\par}
  \vspace{0.6cm}
  {\color{white}\normalsize 2025 --- 2026\par}
  \vfill
\end{titlepage}
\nopagecolor

\tableofcontents
\newpage

% ══════════════════════════════════════════════════════════════════
\section{Introducción}

Suponemos conocidas las medidas de los ángulos y su relación con los
arcos que provienen de la geometría elemental.

La medida $\alpha$ de un ángulo $AOB$ es la \textit{cantidad de
rotación} que efectúa el lado $OA$ al girar en torno a $O$ hasta
coincidir con $OB$. Es positiva en sentido antihorario y negativa en
sentido horario.

\subsection{Sistema Sexagesimal}
\[
  1^\circ = \frac{1}{360}\ \text{de circunferencia completa}
\]

\subsection{Sistema Radiánico}
\[
  1\ \text{rad} = \frac{\text{arco}}{\text{radio}}
  \qquad\Rightarrow\qquad \text{arco} = r
\]

\newpage

% ══════════════════════════════════════════════════════════════════
\section{Razones Trigonométricas}

\begin{tcolorbox}[formulabox,
  title={Definiciones en triángulo rectángulo con ángulo $\alpha$}]
\begin{align*}
\sin\alpha &= \frac{\text{cateto opuesto}}{\text{hipotenusa}}
           = \frac{PQ}{OP} &
\cos\alpha &= \frac{\text{cateto adyacente}}{\text{hipotenusa}}
           = \frac{OQ}{OP} \\[6pt]
\tan\alpha &= \frac{\text{cateto opuesto}}{\text{cateto adyacente}}
           = \frac{PQ}{OQ} &
\cot\alpha &= \frac{\text{cateto adyacente}}{\text{cateto opuesto}}
           = \frac{OQ}{PQ} \\[6pt]
\sec\alpha &= \frac{\text{hipotenusa}}{\text{cateto adyacente}}
           = \frac{OP}{OQ} &
\csc\alpha &= \frac{\text{hipotenusa}}{\text{cateto opuesto}}
           = \frac{OP}{PQ}
\end{align*}
\end{tcolorbox}

% ══════════════════════════════════════════════════════════════════
\section{Identidades Trigonométricas Básicas}

\subsection{Identidades Pitagóricas}

\begin{tcolorbox}[formulabox, title={Identidades Pitagóricas}]
\begin{align*}
\sin^2\alpha + \cos^2\alpha &= 1 \tag{1} \\
1 + \tan^2\alpha            &= \sec^2\alpha \tag{2} \\
1 + \cot^2\alpha            &= \csc^2\alpha \tag{3}
\end{align*}
\end{tcolorbox}

\subsection{Identidades Recíprocas y de Cociente}

\begin{tcolorbox}[formulabox]
\begin{align*}
\sin\alpha\cdot\csc\alpha &= 1 &
\cos\alpha\cdot\sec\alpha &= 1 &
\tan\alpha\cdot\cot\alpha &= 1 \\[4pt]
\tan\alpha &= \dfrac{\sin\alpha}{\cos\alpha} &
\cot\alpha &= \dfrac{\cos\alpha}{\sin\alpha}
\end{align*}
\end{tcolorbox}

\begin{tcolorbox}[notabox]
\textbf{Nota:} $\sin^2\alpha = (\sin\alpha)^2 \neq \sin(\alpha^2)$
\end{tcolorbox}

% ══════════════════════════════════════════════════════════════════
\section{Signos según Cuadrantes}

\begin{center}
\begin{tabular}{>{\centering\bfseries\color{azulOscuro}}m{2cm}
                >{\centering}m{3.4cm}
                >{\centering}m{3.4cm}
                >{\centering\arraybackslash}m{3.6cm}}
\toprule
\rowcolor{azulOscuro}
\color{white}\textbf{Cuadrante} &
\color{white}\textbf{Seno / Cosecante} &
\color{white}\textbf{Coseno / Secante} &
\color{white}\textbf{Tangente / Cotangente} \\
\midrule
\rowcolor{grisClaro}
I   & $+$ & $+$ & $+$ \\
II  & $+$ & $-$ & $-$ \\
\rowcolor{grisClaro}
III & $-$ & $-$ & $+$ \\
IV  & $-$ & $+$ & $-$ \\
\bottomrule
\end{tabular}
\end{center}

% ══════════════════════════════════════════════════════════════════
\section{Fórmulas de Reducción y Cofunciones}

Para $0 < \alpha < \dfrac{\pi}{2}$:

\begin{tcolorbox}[formulabox]
\begin{align*}
\sin\!\left(\tfrac{\pi}{2}-\alpha\right) &= \cos\alpha, &
\cos\!\left(\tfrac{\pi}{2}-\alpha\right) &= \sin\alpha, &
\tan\!\left(\tfrac{\pi}{2}-\alpha\right) &= \cot\alpha \\[4pt]
\sec\!\left(\tfrac{\pi}{2}-\alpha\right) &= \csc\alpha, &
\csc\!\left(\tfrac{\pi}{2}-\alpha\right) &= \sec\alpha
\end{align*}
\end{tcolorbox}

\subsection{Reglas Generales de Reducción}

\begin{tcolorbox}[formulabox]
\begin{align*}
\sin(\pi\pm\alpha) &= \mp\sin\alpha &
\cos(\pi\pm\alpha) &= \mp\cos\alpha &
\tan(\pi\pm\alpha) &= \pm\tan\alpha \\[4pt]
\sin(2\pi\pm\alpha) &= \pm\sin\alpha &
\cos(2\pi\pm\alpha) &= \pm\cos\alpha
\end{align*}
\end{tcolorbox}

\subsection{Ángulos Negativos}

\begin{tcolorbox}[formulabox]
\begin{align*}
\sin(-\theta) &= -\sin\theta &
\cos(-\theta) &= \cos\theta &
\tan(-\theta) &= -\tan\theta \\[4pt]
\csc(-\theta) &= -\csc\theta &
\sec(-\theta) &= \sec\theta &
\cot(-\theta) &= -\cot\theta
\end{align*}
\end{tcolorbox}

\newpage

% ══════════════════════════════════════════════════════════════════
\section{Identidades de Doble, Triple y Semángulo}

\subsection{Doble Ángulo}

\begin{tcolorbox}[formulabox, title={Doble ángulo}]
\begin{align*}
\sin 2\alpha &= 2\sin\alpha\cos\alpha \\[4pt]
\cos 2\alpha &= \cos^2\alpha - \sin^2\alpha
             = 2\cos^2\alpha - 1
             = 1 - 2\sin^2\alpha \\[4pt]
\tan 2\alpha &= \frac{2\tan\alpha}{1-\tan^2\alpha}
\end{align*}
\end{tcolorbox}

\subsection{Triple Ángulo}

\begin{tcolorbox}[formulabox, title={Triple ángulo}]
\begin{align*}
\sin 3\alpha &= 3\sin\alpha - 4\sin^3\alpha \\[4pt]
\cos 3\alpha &= 4\cos^3\alpha - 3\cos\alpha
\end{align*}
\end{tcolorbox}

\subsection{Semángulo}

\begin{tcolorbox}[formulabox, title={Mitad de ángulo}]
\begin{align*}
\sin^2\frac{\alpha}{2} &= \frac{1-\cos\alpha}{2} \\[8pt]
\cos^2\frac{\alpha}{2} &= \frac{1+\cos\alpha}{2} \\[8pt]
\tan\frac{\alpha}{2}   &= \frac{\sin\alpha}{1+\cos\alpha}
                        = \frac{1-\cos\alpha}{\sin\alpha}
                        = \pm\sqrt{\frac{1-\cos\alpha}{1+\cos\alpha}}
\end{align*}
\end{tcolorbox}

% ══════════════════════════════════════════════════════════════════
\section{Identidades de Suma y Diferencia}

\begin{tcolorbox}[formulabox, title={Suma y diferencia de ángulos}]
\begin{align*}
\sin(\alpha\pm\beta) &= \sin\alpha\cos\beta \pm \cos\alpha\sin\beta \\[6pt]
\cos(\alpha\pm\beta) &= \cos\alpha\cos\beta \mp \sin\alpha\sin\beta \\[6pt]
\tan(\alpha\pm\beta) &= \frac{\tan\alpha\pm\tan\beta}
                            {1\mp\tan\alpha\tan\beta}
\end{align*}
\end{tcolorbox}

% ══════════════════════════════════════════════════════════════════
\section{Fórmulas Producto a Suma}

\subsection{Producto hacia Suma}

\begin{tcolorbox}[formulabox]
\begin{align*}
\sin\alpha\cos\beta &= \tfrac{1}{2}\bigl[\sin(\alpha-\beta)+\sin(\alpha+\beta)\bigr] \\[6pt]
\cos\alpha\cos\beta &= \tfrac{1}{2}\bigl[\cos(\alpha-\beta)+\cos(\alpha+\beta)\bigr] \\[6pt]
\sin\alpha\sin\beta &= \tfrac{1}{2}\bigl[\cos(\alpha-\beta)-\cos(\alpha+\beta)\bigr]
\end{align*}
\end{tcolorbox}

\subsection{Suma hacia Producto}

\begin{tcolorbox}[formulabox]
\begin{align*}
\sin\alpha+\sin\beta &= 2\sin\!\left(\frac{\alpha+\beta}{2}\right)
                         \cos\!\left(\frac{\alpha-\beta}{2}\right) \\[8pt]
\sin\alpha-\sin\beta &= 2\cos\!\left(\frac{\alpha+\beta}{2}\right)
                         \sin\!\left(\frac{\alpha-\beta}{2}\right) \\[8pt]
\cos\alpha+\cos\beta &= 2\cos\!\left(\frac{\alpha+\beta}{2}\right)
                         \cos\!\left(\frac{\alpha-\beta}{2}\right) \\[8pt]
\cos\alpha-\cos\beta &= -2\sin\!\left(\frac{\alpha+\beta}{2}\right)
                          \sin\!\left(\frac{\alpha-\beta}{2}\right)
\end{align*}
\end{tcolorbox}

\newpage

% ══════════════════════════════════════════════════════════════════
\section{Teoremas del Seno y del Coseno}

\subsection{Teorema del Seno}

\begin{tcolorbox}[formulabox,
  title={Teorema del Seno --- triángulo $ABC$ oblicuángulo}]
\[
  \frac{\sin A}{a} = \frac{\sin B}{b} = \frac{\sin C}{c}
\]
\end{tcolorbox}

\subsection{Teorema del Coseno}

\begin{tcolorbox}[formulabox, title={Teorema del Coseno}]
\begin{align*}
a^2 &= b^2 + c^2 - 2bc\cos A \\[4pt]
b^2 &= a^2 + c^2 - 2ac\cos B \\[4pt]
c^2 &= a^2 + b^2 - 2ab\cos C
\end{align*}
\end{tcolorbox}

\subsection{Área del Triángulo}

\begin{tcolorbox}[formulabox]
\[
  \text{Área} = \frac{1}{2}\,bc\sin A
              = \frac{abc}{4R}
  \qquad R = \text{radio circunscrito}
\]
\end{tcolorbox}

% ══════════════════════════════════════════════════════════════════
\section{Valores Exactos de Ángulos Notables}

\begin{center}
\begin{tabular}{>{\centering\bfseries}m{2cm}
                >{\centering}m{3cm}
                >{\centering}m{3cm}
                >{\centering\arraybackslash}m{3cm}}
\toprule
\rowcolor{azulOscuro}
\color{white}\textbf{Ángulo} &
\color{white}$\boldsymbol{\sin}$ &
\color{white}$\boldsymbol{\cos}$ &
\color{white}$\boldsymbol{\tan}$ \\
\midrule
\rowcolor{grisClaro}
$0^\circ$  & $0$ & $1$ & $0$ \\
$30^\circ$ & $\dfrac{1}{2}$ & $\dfrac{\sqrt{3}}{2}$ &
             $\dfrac{1}{\sqrt{3}}$ \\[10pt]
\rowcolor{grisClaro}
$45^\circ$ & $\dfrac{\sqrt{2}}{2}$ & $\dfrac{\sqrt{2}}{2}$ & $1$ \\[10pt]
$60^\circ$ & $\dfrac{\sqrt{3}}{2}$ & $\dfrac{1}{2}$ & $\sqrt{3}$ \\[6pt]
\rowcolor{grisClaro}
$90^\circ$ & $1$ & $0$ & $\infty$ \\
\bottomrule
\end{tabular}
\end{center}

\subsection{Triángulos Notables}
\begin{itemize}[leftmargin=*, itemsep=2pt]
  \item Triángulo equilátero de lado 2
        $\;\longrightarrow\;$ razones de $30^\circ$ y $60^\circ$
  \item Triángulo rectángulo isósceles de catetos 1
        $\;\longrightarrow\;$ razones de $45^\circ$
\end{itemize}

\newpage

% ══════════════════════════════════════════════════════════════════
\section{Otras Identidades Útiles}

\begin{tcolorbox}[formulabox, title={Demostraciones frecuentes}]
\begin{align*}
\sec\theta  - \cos\theta  &= \tan\theta\sin\theta \\[4pt]
\csc\theta  - \sin\theta  &= \cot\theta\cos\theta \\[4pt]
\tan\theta  + \cot\theta  &= \csc\theta\sec\theta \\[4pt]
\sin^2\theta - \cos^2\theta &= -\cos 2\theta \\[4pt]
\sin^3\theta + \cos^3\theta
  &= \bigl(\sin\theta+\cos\theta\bigr)
     \bigl(1-\sin\theta\cos\theta\bigr)
\end{align*}
\end{tcolorbox}

% ══════════════════════════════════════════════════════════════════
\section{Funciones Trigonométricas Asociadas --- Ondas}

Forma general:
\[
  y = A\sin(Bx+C)+D \qquad\text{o}\qquad y = A\cos(Bx+C)+D
\]

\begin{tcolorbox}[formulabox]
\begin{itemize}[leftmargin=*, itemsep=4pt]
  \item \textbf{Amplitud:} $|A|$
  \item \textbf{Período:} $\dfrac{2\pi}{|B|}$
  \item \textbf{Frecuencia:} $\dfrac{|B|}{2\pi}$
  \item \textbf{Desfase:} $-\dfrac{C}{B}$
  \item \textbf{Desplazamiento vertical:} $D$
\end{itemize}
\end{tcolorbox}

\[
  A = \frac{\text{Valor máximo} - \text{Valor mínimo}}{2}
\]

% ══════════════════════════════════════════════════════════════════
\subsection{Gráfico de la Función Seno}

La función seno base $y = \sin(x)$ tiene amplitud $1$, período $2\pi$,
desfase $0$ y desplazamiento vertical $0$.

\begin{center}
\begin{tikzpicture}
\begin{axis}[
  width=12cm, height=4.8cm,
  axis lines=center,
  xlabel={$x$}, ylabel={$y$},
  xlabel style={right},
  ylabel style={above},
  xtick={-6.2832,-3.1416,0,1.5708,3.1416,6.2832},
  xticklabels={$-2\pi$,$-\pi$,$0$,$\tfrac{\pi}{2}$,$\pi$,$2\pi$},
  ytick={-2,-1,1,2},
  xmin=-7, xmax=7,
  ymin=-2.5, ymax=2.5,
  grid=major,
  grid style={line width=0.4pt, draw=gray!40},
  tick label style={font=\small},
  domain=-6.5:6.5,
  samples=80,
  legend pos=north east,
  legend style={font=\small, fill=grisClaro, draw=azulMedio}
]
  \addplot[color=azulOscuro, thick, smooth] {sin(deg(x))};
  \addlegendentry{$y = \sin(x)$}
  \addplot[color=rojoGrafico, thick, dashed, smooth] {2*sin(deg(x))};
  \addlegendentry{$y = 2\sin(x)$ \quad amplitud $= 2$}
\end{axis}
\end{tikzpicture}
\end{center}

% ══════════════════════════════════════════════════════════════════
\subsection{Gráfico de la Función Coseno}

La función coseno base $y = \cos(x)$ tiene amplitud $1$, período $2\pi$,
desfase $0$ y desplazamiento vertical $0$. A diferencia del seno,
comienza en su valor máximo cuando $x = 0$.

\begin{center}
\begin{tikzpicture}
\begin{axis}[
  width=12cm, height=4.8cm,
  axis lines=center,
  xlabel={$x$}, ylabel={$y$},
  xlabel style={right},
  ylabel style={above},
  xtick={-6.2832,-3.1416,0,1.5708,3.1416,6.2832},
  xticklabels={$-2\pi$,$-\pi$,$0$,$\tfrac{\pi}{2}$,$\pi$,$2\pi$},
  ytick={-2,-1,1,2},
  xmin=-7, xmax=7,
  ymin=-2.5, ymax=2.5,
  grid=major,
  grid style={line width=0.4pt, draw=gray!40},
  tick label style={font=\small},
  domain=-6.5:6.5, samples=80,
  legend pos=north east,
  legend style={font=\small, fill=grisClaro, draw=azulMedio}
]
  \addplot[color=verdeGrafico, thick, smooth] {cos(deg(x))};
  \addlegendentry{$y = \cos(x)$}
  \addplot[color=rojoGrafico, thick, dashed, smooth] {1.5*cos(deg(x))};
  \addlegendentry{$y = 1.5\cos(x)$ \quad amplitud $= 1.5$}
  \addplot[color=azulOscuro, thick, dashdotted, smooth] {cos(deg(x))+1};
  \addlegendentry{$y = \cos(x)+1$ \quad $D = 1$}
\end{axis}
\end{tikzpicture}
\end{center}

% ══════════════════════════════════════════════════════════════════
\subsection{Efecto del Período y el Desfase}

\begin{center}
\begin{tikzpicture}
\begin{axis}[
  width=12cm, height=4.8cm,
  axis lines=center,
  xlabel={$x$}, ylabel={$y$},
  xlabel style={right},
  ylabel style={above},
  xtick={-6.2832,-3.1416,0,1.5708,3.1416,6.2832},
  xticklabels={$-2\pi$,$-\pi$,$0$,$\tfrac{\pi}{2}$,$\pi$,$2\pi$},
  ytick={-1,1},
  xmin=-7, xmax=7,
  ymin=-1.6, ymax=1.6,
  grid=major,
  grid style={line width=0.4pt, draw=gray!40},
  tick label style={font=\small},
  domain=-6.5:6.5, samples=80,
  legend pos=north east,
  legend style={font=\small, fill=grisClaro, draw=azulMedio}
]
  \addplot[color=verdeGrafico, thick, smooth] {cos(deg(x))};
  \addlegendentry{$y = \cos(x)$ \quad $T = 2\pi$}
  \addplot[color=azulOscuro, thick, dashed, smooth] {cos(deg(2*x))};
  \addlegendentry{$y = \cos(2x)$ \quad $T = \pi$}
  \addplot[color=rojoGrafico, thick, dashdotted, smooth] {cos(deg(x-1.0472))};
  \addlegendentry{$y = \cos(x - \pi/3)$ \quad desfase $= \pi/3$}
\end{axis}
\end{tikzpicture}
\end{center}

\begin{tcolorbox}[notabox]
\textbf{Observación clave:} $\cos(x) = \sin\!\left(x + \dfrac{\pi}{2}\right)$.
El coseno es el seno desplazado $\tfrac{\pi}{2}$ a la \textit{izquierda}.
\end{tcolorbox}

% ══════════════════════════════════════════════════════════════════
%  NUEVA SUBSECCIÓN 12.4 — Tangente, Cotangente, Secante, Cosecante
% ══════════════════════════════════════════════════════════════════
\subsection{Tangente, Cotangente, Secante y Cosecante}

Forma general para cada función:
\[
  y = A\,f(Bx + C) + D
\]
donde $f$ representa $\tan$, $\cot$, $\sec$ o $\csc$.

% ── TANGENTE ─────────────────────────────────────────────────────
\subsubsection*{$y = A\tan(Bx+C)+D$}

\begin{itemize}[leftmargin=*, itemsep=3pt]
  \item \textbf{Amplitud (factor de escala vertical):} $|A|$
        \quad{\small(estira/comprime la curva; no existe máximo ni mínimo acotado)}
  \item \textbf{Período:} $\dfrac{\pi}{|B|}$
  \item \textbf{Desfase:} $-\dfrac{C}{B}$
        \quad{\small(positivo $\Rightarrow$ desplaza a la derecha)}
  \item \textbf{Desplazamiento vertical:} $D$
  \item \textbf{Asíntotas verticales:}
        $\displaystyle x = \frac{1}{B}\!\left(\frac{\pi}{2}+n\pi-C\right)$,\;
        $n\in\mathbb{Z}$
  \item \textbf{Ejemplos graficados:}\\[2pt]
        $y=\tan(x)$: $A=1$, $T=\pi$, desfase $=0$\\
        $y=2\tan(x)$: $A=2$, $T=\pi$, desfase $=0$\\
        $y=\tan(2x)$: $A=1$, $T=\tfrac{\pi}{2}$, desfase $=0$\\
        $y=\tan\!\left(x-\tfrac{\pi}{4}\right)$:
          $A=1$, $T=\pi$, desfase $=\tfrac{\pi}{4}$
\end{itemize}

\begin{center}
\begin{tikzpicture}
\begin{axis}[
  width=12cm, height=5cm,
  axis lines=center,
  xlabel={$x$}, ylabel={$y$},
  xlabel style={right}, ylabel style={above},
  xtick={-4.7124,-3.1416,-1.5708,0,1.5708,3.1416,4.7124},
  xticklabels={$-\frac{3\pi}{2}$,$-\pi$,$-\frac{\pi}{2}$,$0$,
               $\frac{\pi}{2}$,$\pi$,$\frac{3\pi}{2}$},
  ytick={-3,-2,-1,1,2,3},
  xmin=-5.2, xmax=5.2,
  ymin=-3.5, ymax=3.5,
  restrict y to domain=-3.5:3.5,
  grid=major,
  grid style={line width=0.4pt, draw=gray!40},
  tick label style={font=\small},
  domain=-5.1:5.1, samples=200,
  legend pos=north east,
  legend style={font=\small, fill=grisClaro, draw=azulMedio}
]
  % Asíntotas verticales
  \addplot[gray!50, dashed, thick] coordinates {(-1.5708,-3.5)(-1.5708,3.5)};
  \addplot[gray!50, dashed, thick] coordinates {( 1.5708,-3.5)( 1.5708,3.5)};
  \addplot[gray!50, dashed, thick] coordinates {(-4.7124,-3.5)(-4.7124,3.5)};
  \addplot[gray!50, dashed, thick] coordinates {( 4.7124,-3.5)( 4.7124,3.5)};
  % Curvas
  \addplot[azulOscuro, thick, smooth]               {tan(deg(x))};
  \addlegendentry{$y=\tan(x)$\quad $T=\pi$}
  \addplot[rojoGrafico, thick, dashed, smooth]      {2*tan(deg(x))};
  \addlegendentry{$y=2\tan(x)$\quad $A=2$}
  \addplot[verdeGrafico, thick, dashdotted, smooth] {tan(deg(2*x))};
  \addlegendentry{$y=\tan(2x)$\quad $T=\pi/2$}
  \addplot[dorado, thick, dotted, smooth]           {tan(deg(x-0.7854))};
  \addlegendentry{$y=\tan(x-\pi/4)$\quad desfase$\,=\pi/4$}
\end{axis}
\end{tikzpicture}
\end{center}

\newpage

% ── COTANGENTE ───────────────────────────────────────────────────
\subsubsection*{$y = A\cot(Bx+C)+D$}

\begin{itemize}[leftmargin=*, itemsep=3pt]
  \item \textbf{Amplitud (factor de escala vertical):} $|A|$
        \quad{\small(función \textbf{decreciente} en cada período; sin máximo/mínimo acotado)}
  \item \textbf{Período:} $\dfrac{\pi}{|B|}$
  \item \textbf{Desfase:} $-\dfrac{C}{B}$
  \item \textbf{Desplazamiento vertical:} $D$
  \item \textbf{Asíntotas verticales:}
        $\displaystyle x = \frac{n\pi - C}{B}$,\; $n\in\mathbb{Z}$
  \item \textbf{Diferencia clave:} la cotangente es decreciente en cada intervalo,
        mientras que la tangente es creciente.
  \item \textbf{Ejemplos graficados:}\\[2pt]
        $y=\cot(x)$: $A=1$, $T=\pi$, desfase $=0$\\
        $y=2\cot(x)$: $A=2$, $T=\pi$, desfase $=0$\\
        $y=\cot(2x)$: $A=1$, $T=\tfrac{\pi}{2}$, desfase $=0$\\
        $y=\cot\!\left(x+\tfrac{\pi}{4}\right)$:
          $A=1$, $T=\pi$, desfase $=-\tfrac{\pi}{4}$
\end{itemize}

\begin{center}
\begin{tikzpicture}
\begin{axis}[
  width=12cm, height=5cm,
  axis lines=center,
  xlabel={$x$}, ylabel={$y$},
  xlabel style={right}, ylabel style={above},
  xtick={-3.1416,-1.5708,0,1.5708,3.1416,4.7124},
  xticklabels={$-\pi$,$-\frac{\pi}{2}$,$0$,
               $\frac{\pi}{2}$,$\pi$,$\frac{3\pi}{2}$},
  ytick={-3,-2,-1,1,2,3},
  xmin=-3.8, xmax=5.2,
  ymin=-3.5, ymax=3.5,
  restrict y to domain=-3.5:3.5,
  grid=major,
  grid style={line width=0.4pt, draw=gray!40},
  tick label style={font=\small},
  domain=-3.7:5.1, samples=200,
  legend pos=north east,
  legend style={font=\small, fill=grisClaro, draw=azulMedio}
]
  % Asíntotas verticales
  \addplot[gray!50, dashed, thick] coordinates {( 0,    -3.5)( 0,    3.5)};
  \addplot[gray!50, dashed, thick] coordinates {( 3.1416,-3.5)( 3.1416,3.5)};
  \addplot[gray!50, dashed, thick] coordinates {(-3.1416,-3.5)(-3.1416,3.5)};
  % Curvas
  \addplot[azulOscuro, thick, smooth]
    {cos(deg(x))/sin(deg(x))};
  \addlegendentry{$y=\cot(x)$\quad $T=\pi$}
  \addplot[rojoGrafico, thick, dashed, smooth]
    {2*cos(deg(x))/sin(deg(x))};
  \addlegendentry{$y=2\cot(x)$\quad $A=2$}
  \addplot[verdeGrafico, thick, dashdotted, smooth]
    {cos(deg(2*x))/sin(deg(2*x))};
  \addlegendentry{$y=\cot(2x)$\quad $T=\pi/2$}
  \addplot[dorado, thick, dotted, smooth]
    {cos(deg(x+0.7854))/sin(deg(x+0.7854))};
  \addlegendentry{$y=\cot(x+\pi/4)$\quad desfase$\,=-\pi/4$}
\end{axis}
\end{tikzpicture}
\end{center}

\newpage

% ── SECANTE ──────────────────────────────────────────────────────
\subsubsection*{$y = A\sec(Bx+C)+D$}

\begin{itemize}[leftmargin=*, itemsep=3pt]
  \item \textbf{Amplitud (factor de escala vertical):} $|A|$
        \quad{\small(los vértices de las ramas se ubican en $\pm|A|+D$)}
  \item \textbf{Período:} $\dfrac{2\pi}{|B|}$
  \item \textbf{Desfase:} $-\dfrac{C}{B}$
  \item \textbf{Desplazamiento vertical:} $D$
  \item \textbf{Asíntotas verticales:}
        $\displaystyle x = \frac{1}{B}\!\left(\frac{\pi}{2}+n\pi-C\right)$,\;
        $n\in\mathbb{Z}$ \quad{\small(donde $\cos = 0$)}
  \item \textbf{Relación:} $\sec(x)=\dfrac{1}{\cos(x)}$;\;
        la curva \textit{toca} al coseno en sus máximos y mínimos.
  \item \textbf{Ejemplos graficados:}\\[2pt]
        $y=\sec(x)$: $A=1$, $T=2\pi$, desfase $=0$\\
        $y=1{,}5\sec(x)$: $A=1{,}5$, $T=2\pi$, desfase $=0$\\
        $y=\sec(2x)$: $A=1$, $T=\pi$, desfase $=0$\\
        $y=\sec\!\left(x-\tfrac{\pi}{3}\right)$:
          $A=1$, $T=2\pi$, desfase $=\tfrac{\pi}{3}$
\end{itemize}

\begin{center}
\begin{tikzpicture}
\begin{axis}[
  width=12cm, height=5cm,
  axis lines=center,
  xlabel={$x$}, ylabel={$y$},
  xlabel style={right}, ylabel style={above},
  xtick={-4.7124,-3.1416,-1.5708,0,1.5708,3.1416,4.7124},
  xticklabels={$-\frac{3\pi}{2}$,$-\pi$,$-\frac{\pi}{2}$,$0$,
               $\frac{\pi}{2}$,$\pi$,$\frac{3\pi}{2}$},
  ytick={-3,-2,-1,1,2,3},
  xmin=-5.2, xmax=5.2,
  ymin=-3.5, ymax=3.5,
  restrict y to domain=-3.5:3.5,
  grid=major,
  grid style={line width=0.4pt, draw=gray!40},
  tick label style={font=\small},
  domain=-5.1:5.1, samples=200,
  legend pos=north east,
  legend style={font=\small, fill=grisClaro, draw=azulMedio}
]
  % Asíntotas verticales
  \addplot[gray!50, dashed, thick] coordinates {(-1.5708,-3.5)(-1.5708,3.5)};
  \addplot[gray!50, dashed, thick] coordinates {( 1.5708,-3.5)( 1.5708,3.5)};
  \addplot[gray!50, dashed, thick] coordinates {(-4.7124,-3.5)(-4.7124,3.5)};
  \addplot[gray!50, dashed, thick] coordinates {( 4.7124,-3.5)( 4.7124,3.5)};
  % cos(x) de referencia
  \addplot[gray!55, thin, dotted, smooth] {cos(deg(x))};
  % Curvas
  \addplot[azulOscuro, thick, smooth]               {1/cos(deg(x))};
  \addlegendentry{$y=\sec(x)$\quad $T=2\pi$}
  \addplot[rojoGrafico, thick, dashed, smooth]      {1.5/cos(deg(x))};
  \addlegendentry{$y=1{,}5\sec(x)$\quad $A=1{,}5$}
  \addplot[verdeGrafico, thick, dashdotted, smooth] {1/cos(deg(2*x))};
  \addlegendentry{$y=\sec(2x)$\quad $T=\pi$}
  \addplot[dorado, thick, dotted, smooth]           {1/cos(deg(x-1.0472))};
  \addlegendentry{$y=\sec(x-\pi/3)$\quad desfase$\,=\pi/3$}
\end{axis}
\end{tikzpicture}
\end{center}

\begin{tcolorbox}[notabox]
\textbf{Nota:} La línea punteada gris es $y=\cos(x)$ como referencia.
La secante toca los vértices del coseno y se extiende hacia $\pm\infty$
en sus asíntotas.
\end{tcolorbox}

\newpage

% ── COSECANTE ────────────────────────────────────────────────────
\subsubsection*{$y = A\csc(Bx+C)+D$}

\begin{itemize}[leftmargin=*, itemsep=3pt]
  \item \textbf{Amplitud (factor de escala vertical):} $|A|$
        \quad{\small(los vértices de las ramas se ubican en $\pm|A|+D$)}
  \item \textbf{Período:} $\dfrac{2\pi}{|B|}$
  \item \textbf{Desfase:} $-\dfrac{C}{B}$
  \item \textbf{Desplazamiento vertical:} $D$
  \item \textbf{Asíntotas verticales:}
        $\displaystyle x = \frac{n\pi - C}{B}$,\; $n\in\mathbb{Z}$
        \quad{\small(donde $\sin = 0$)}
  \item \textbf{Relación:} $\csc(x)=\dfrac{1}{\sin(x)}$;\;
        la curva \textit{toca} al seno en sus máximos y mínimos.
  \item \textbf{Ejemplos graficados:}\\[2pt]
        $y=\csc(x)$: $A=1$, $T=2\pi$, desfase $=0$\\
        $y=1{,}5\csc(x)$: $A=1{,}5$, $T=2\pi$, desfase $=0$\\
        $y=\csc(2x)$: $A=1$, $T=\pi$, desfase $=0$\\
        $y=\csc\!\left(x+\tfrac{\pi}{6}\right)$:
          $A=1$, $T=2\pi$, desfase $=-\tfrac{\pi}{6}$
\end{itemize}

\begin{center}
\begin{tikzpicture}
\begin{axis}[
  width=12cm, height=5cm,
  axis lines=center,
  xlabel={$x$}, ylabel={$y$},
  xlabel style={right}, ylabel style={above},
  xtick={-3.1416,-1.5708,0,1.5708,3.1416,4.7124,6.2832},
  xticklabels={$-\pi$,$-\frac{\pi}{2}$,$0$,
               $\frac{\pi}{2}$,$\pi$,$\frac{3\pi}{2}$,$2\pi$},
  ytick={-3,-2,-1,1,2,3},
  xmin=-3.8, xmax=6.8,
  ymin=-3.5, ymax=3.5,
  restrict y to domain=-3.5:3.5,
  grid=major,
  grid style={line width=0.4pt, draw=gray!40},
  tick label style={font=\small},
  domain=-3.7:6.7, samples=200,
  legend pos=north east,
  legend style={font=\small, fill=grisClaro, draw=azulMedio}
]
  % Asíntotas verticales
  \addplot[gray!50, dashed, thick] coordinates {( 0,    -3.5)( 0,    3.5)};
  \addplot[gray!50, dashed, thick] coordinates {( 3.1416,-3.5)( 3.1416,3.5)};
  \addplot[gray!50, dashed, thick] coordinates {( 6.2832,-3.5)( 6.2832,3.5)};
  \addplot[gray!50, dashed, thick] coordinates {(-3.1416,-3.5)(-3.1416,3.5)};
  % sin(x) de referencia
  \addplot[gray!55, thin, dotted, smooth] {sin(deg(x))};
  % Curvas
  \addplot[azulOscuro, thick, smooth]               {1/sin(deg(x))};
  \addlegendentry{$y=\csc(x)$\quad $T=2\pi$}
  \addplot[rojoGrafico, thick, dashed, smooth]      {1.5/sin(deg(x))};
  \addlegendentry{$y=1{,}5\csc(x)$\quad $A=1{,}5$}
  \addplot[verdeGrafico, thick, dashdotted, smooth] {1/sin(deg(2*x))};
  \addlegendentry{$y=\csc(2x)$\quad $T=\pi$}
  \addplot[dorado, thick, dotted, smooth]           {1/sin(deg(x+0.5236))};
  \addlegendentry{$y=\csc(x+\pi/6)$\quad desfase$\,=-\pi/6$}
\end{axis}
\end{tikzpicture}
\end{center}

\begin{tcolorbox}[notabox]
\textbf{Nota:} La línea punteada gris es $y=\sin(x)$ como referencia.
La cosecante toca los vértices del seno y se extiende hacia $\pm\infty$
en los ceros del seno.
\end{tcolorbox}

\newpage

% ══════════════════════════════════════════════════════════════════
\section{Identidades para Demostrar}

\begin{tcolorbox}[ejercbox,
  title={\color{dorado}\bfseries Ejercicios en clases}]
\begin{align*}
a)\quad
  &\frac{\csc\alpha}{1+\csc\alpha}
   - \frac{1}{\sin\alpha-1} = 2\sec^2\alpha \\[8pt]
b)\quad
  &\frac{\tan\alpha-\cot\alpha}{1-2\cos^2\alpha}
   = \tan\alpha+\cot\alpha \\[8pt]
c)\quad
  &\frac{\sin\alpha}{1+\cos\alpha}+\cot\alpha = \csc\alpha \\[8pt]
d)\quad
  &\frac{\sec\alpha+1}{\sec\alpha-1}
   -\frac{\sec\alpha-1}{\sec\alpha+1}
   -4\cot^2\alpha
   = \frac{4}{1+\sec\alpha} \\[8pt]
e)\quad
  &\frac{1+\tan^2\alpha}{1+\cot^2\alpha}
   = \left(\frac{1-\tan\alpha}{1-\cot\alpha}\right)^{\!2} \\[8pt]
f)\quad
  &\frac{1-\sin\alpha\cos\alpha}{\cos\alpha\,(\sec\alpha-\csc\alpha)}
   \cdot
   \frac{\sin^2\alpha-\cos^2\alpha}{\sin^3\alpha+\cos^3\alpha}
   = \sin\alpha
\end{align*}
\end{tcolorbox}

% ══════════════════════════════════════════════════════════════════
\section{Trabajo Autónomo --- Identidades Adicionales}

\begin{tcolorbox}[ejercbox,
  title={\color{dorado}\bfseries Trabajo autónomo}]
\begin{align*}
g)\quad
  &2\sec^2\alpha - \sec^4\alpha - 2\cot^2\alpha + \cot^4\alpha
   = \cot^4\alpha - \tan^4\alpha \\[8pt]
h)\quad
  &\frac{1}{1+\sin^2\alpha}+\frac{1}{1+\cos^2\alpha}
   = \csc^4\alpha - \cot^2\alpha\,(2+\cot^2\alpha) \\[8pt]
i)\quad
  &\frac{\sin^6\alpha-\cos^6\alpha}{\sin^2\alpha-\cos^2\alpha}
   +\sin^2\alpha\cos^2\alpha = 1 \\[8pt]
j)\quad
  &\frac{\tan^3\alpha}{1+\tan^2\alpha}
   +\frac{\cot^3\alpha}{1+\cot^2\alpha}
   = \sec\alpha\,\csc\alpha - 2\sin\alpha\cos\alpha
\end{align*}
\end{tcolorbox}

\vspace{4pt}

\begin{tcolorbox}[ejercbox, title={\color{dorado}\bfseries Identidades adicionales (65--73)}]
\begin{align*}
k)\quad &\sec\theta - \cos\theta = \tan\theta\sin\theta \\[6pt]
l)\quad &\csc\theta - \sin\theta = \cot\theta\cos\theta \\[6pt]
m)\quad &\frac{\sin\theta + \cos\theta}{\sin\theta} = 1 + \cot\theta \\[8pt]
n)\quad &\frac{\sin\theta + \cos\theta}{\cos\theta} = 1 + \tan\theta \\[8pt]
\~{n})\quad &\bigl(\cot\theta + \csc\theta\bigr)\bigl(\tan\theta - \sin\theta\bigr)
          = \sec\theta - \cos\theta \\[6pt]
o)\quad &\cot\theta + \tan\theta = \csc\theta\sec\theta \\[6pt]
p)\quad &\sec^2 3\theta\,\csc^2 3\theta = \sec^2 3\theta + \csc^2 3\theta \\[6pt]
q)\quad &\frac{1 + \cos^2 3\theta}{\sin^2 3\theta} = 2\csc^2 3\theta - 1 \\[8pt]
r)\quad &\log\csc\theta = -\log\sin\theta
\end{align*}
\end{tcolorbox}

\newpage

% ══════════════════════════════════════════════════════════════════
\section{Ejercicios de Aplicación}

\begin{tcolorbox}[ejercbox,
  title={\color{dorado}\bfseries Problemas propuestos}]

\begin{enumerate}[leftmargin=*, itemsep=10pt]

  \item \textbf{Dos observadores y un avión.}
  Dos observadores en tierra, separados $d = 1{,}5\ \text{km}$ entre sí,
  observan simultáneamente un avión que vuela horizontalmente entre ellos
  en el mismo plano vertical. El primer observador mide un ángulo de
  elevación de $\alpha = 55{,}3^\circ$ y el segundo uno de $\beta = 72{,}1^\circ$.
  Determine: \textbf{(a)} la altura de vuelo del avión, y
  \textbf{(b)} la velocidad del avión si, al cabo de $t = 10\ \text{s}$
  desde la primera medición, el avión se encuentra exactamente sobre
  el segundo observador.

  \item \textbf{Altura de una torre.}
  Desde un punto $P$ en el suelo se mide un ángulo de elevación de
  $\theta_1 = 38{,}2^\circ$ hasta la cima de una torre vertical.
  Al avanzar $d = 50\ \text{m}$ en línea recta hacia la base de la torre,
  el nuevo ángulo de elevación es $\theta_2 = 61{,}7^\circ$.
  Calcule la altura de la torre.

  \item \textbf{Ancho de un río.}
  Un topógrafo desea medir el ancho de un río. Desde el punto $A$,
  en una orilla, mide un ángulo de $\alpha = 73{,}4^\circ$ hacia una
  piedra $C$ en la orilla opuesta. Luego se desplaza $b = 30\ \text{m}$
  a lo largo de la orilla hasta el punto $B$, desde donde el ángulo
  hacia $C$ es $\beta = 51{,}2^\circ$. Suponiendo que $A$, $B$ y la
  base de la perpendicular desde $C$ son colineales, halle el ancho
  del río.

  \item \textbf{Altura de un acantilado (forma general).}
  Desde el pie de un acantilado vertical, en un punto $P$ del llano,
  se mide el ángulo de elevación $\alpha$ a la cima. Al alejarse una
  distancia $h$ en dirección horizontal hasta el punto $Q$, el ángulo
  de elevación disminuye a $\beta$ ($\beta < \alpha$).
  Demuestre que la altura del acantilado es:
  \[
    H = \frac{h\,\tan\alpha\,\tan\beta}{\tan\alpha - \tan\beta}}
  \]

  \item \textbf{Faro y dos botes.}
  Un faro se encuentra en el punto $F$ de la costa. Dos botes $A$ y $B$
  están en el mar; desde el faro, el ángulo $\angle AFB = 42{,}7^\circ$.
  Se sabe que $FA = 28{,}4\ \text{m}$ y $FB = 20{,}9\ \text{m}$.
  Halle la distancia $AB$ entre los dos botes.

  \item \textbf{Ancho de un río (hombre caminando).}
  Un hombre en la orilla sur de un río observa un árbol en la orilla norte
  con un ángulo de $30^\circ$ al este del norte. Camina $100\ \text{m}$ hacia
  el este a lo largo de la orilla y ahora el árbol se encuentra al norte exacto.
  ¿Cuál es el ancho del río?

  \item \textbf{Acera inclinada.}
  Una acera plana une dos edificios: uno de altura $H_1 = 15\ \text{m}$
  y otro de altura $H_2 = 8\ \text{m}$, cuyas bases están separadas
  horizontalmente $D = 18\ \text{m}$.
  Encuentre el ángulo $\alpha$ de inclinación que forma la acera con
  la horizontal.

\end{enumerate}
\end{tcolorbox}

% ══════════════════════════════════════════════════════════════════
\section{Respuestas}

\begin{tcolorbox}[ejercbox,
  title={\color{dorado}\bfseries Respuestas}]
\begin{enumerate}[leftmargin=*, itemsep=6pt]
  \item Dos observadores y avión:\quad
        altura $1{,}477\ \text{km}$;\quad
        velocidad $= 47{,}7\ \text{m/s}$
  \item Altura de torre: $68{,}3\ \text{m}$
  \item Ancho del río: $27{,}2\ \text{m}$
  \item Altura de acantilado: $H = \dfrac{h\,\tan\alpha\,\tan\beta}{\tan\alpha - \tan\beta}$
  \item Faro y botes: distancia $= 19{,}3\ \text{m}$
  \item Ancho río (hombre caminando): $100\sqrt{3}\ \text{m}$
  \item Acera inclinada: $\alpha \approx 21{,}25^\circ$
\end{enumerate}
\end{tcolorbox}

\vspace{0.8cm}

\begin{tcolorbox}[formulabox, title={Ejercicio avanzado --- Bandas y poleas}]
\[
  \theta + 2\cot\frac{\theta}{2} = \frac{L}{R+r} - \pi
\]
\end{tcolorbox}

\end{document}